% from ega1-10-fr.tex
\oldpage[I]{205}

(\hyperref[0.3.8.2-fr]{0, 3.8.2}). En outre, si $\mathscr{H}$ est un
sous-$\mathscr{O}_{\mathfrak{X}}$-Module cohérent d'un
$\mathscr{O}_{\mathfrak{X}}$-Module cohérent $\mathscr{F}$, pour tout
ouvert $U\subset\mathfrak{X}$, $\Gamma(U,\mathscr{H})$ est un sous-groupe
\emph{fermé} du groupe topologique $\Gamma(U,\mathscr{F})$ car le foncteur
$\Gamma$ étant exact à gauche, $\Gamma(U,\mathscr{H})$ est le noyau de
l'homomorphisme
$\Gamma(U,\mathscr{F})\to\Gamma(U,\mathscr{F}/\mathscr{H})$, qui est
\emph{continu} d'après ce qui précède, puisque $\mathscr{F}/\mathscr{H}$
est cohérent (\hyperref[0.5.3.4-fr]{0, 5.3.4})~; notre assertion résulte
de ce que $\Gamma(U,\mathscr{F}/\mathscr{H})$ est un groupe topologique
séparé.
\end{env}

\begin{proposition}[10.11.7]
\label{I.10.11.7-fr}
Soient $\mathscr{F}$ et $\mathscr{G}$ deux
$\mathscr{O}_{\mathfrak{X}}$-Modules cohérents. On peut définir (avec les
notations de (\hyperref[I.10.11.3-fr]{10.11.3})) des isomorphismes
canoniques fonctoriels de $\mathscr{O}_{\mathfrak{X}}$-Modules
topologiques (\hyperref[I.10.11.6-fr]{10.11.6})
\[
  \label{I.10.11.7.1-fr}
  \mathscr{F}\otimes_{\mathscr{O}_{\mathfrak{X}}}\mathscr{G}
  \xrightarrow{\sim}
  \varprojlim_n
  (\mathscr{F}_n\otimes_{\mathscr{O}_{X_n}}\mathscr{G}_n)
  \tag{10.11.7.1}
\]
\[
  \label{I.10.11.7.2-fr}
  \mathscr{H}\!om_{\mathscr{O}_{\mathfrak{X}}}
  (\mathscr{F},\mathscr{G})
  \xrightarrow{\sim}
  \varprojlim_n
  \mathscr{H}\!om_{\mathscr{O}_{X_n}}(\mathscr{F}_n,\mathscr{G}_n)
  \tag{10.11.7.2}
\]
\end{proposition}

L'existence de l'isomorphisme
(\hyperref[I.10.11.7.1-fr]{10.11.7.1}) résulte de la formule
\[
  \mathscr{F}_n\otimes_{\mathscr{O}_{X_n}}\mathscr{G}_n
  =
  (\mathscr{F}\otimes_{\mathscr{O}_{\mathfrak{X}}}
  \mathscr{O}_{X_n})
  \otimes_{\mathscr{O}_{X_n}}
  (\mathscr{G}\otimes_{\mathscr{O}_{\mathfrak{X}}}
  \mathscr{O}_{X_n})
  =
  (\mathscr{F}\otimes_{\mathscr{O}_{\mathfrak{X}}}\mathscr{G})
  \otimes_{\mathscr{O}_{\mathfrak{X}}}\mathscr{O}_{X_n}
\]
et de (\hyperref[I.10.11.3-fr]{10.11.3}). L'isomorphisme
(\hyperref[I.10.11.7.2-fr]{10.11.7.2}) où les deux membres sont
considérés comme faisceaux de modules sans topologie, résulte de la
définition des sections de
$\mathscr{H}\!om_{\mathscr{O}_{\mathfrak{X}}}
(\mathscr{F},\mathscr{G})$ et
$\mathscr{H}\!om_{\mathscr{O}_{X_n}}(\mathscr{F}_n,\mathscr{G}_n)$ et de
l'existence de l'isomorphisme
(\hyperref[I.10.11.4.1-fr]{10.11.4.1}), appliqué au préschéma induit sur
un ouvert formel affine noethérien arbitraire de $\mathfrak{X}$. Reste à
prouver que l'isomorphisme
(\hyperref[I.10.11.7.2-fr]{10.11.7.2}) est bicontinu au-dessus d'un
ensemble quasi-compact, et on est donc ramené au cas où
$\mathfrak{X}=\operatorname{Spf}(A)$, $A$ étant adique noethérien, d'où
(\hyperref[I.10.10.5-fr]{10.10.5})
$\mathscr{F}=M^\Delta$, $\mathscr{G}=N^\Delta$, $M$, $N$ étant des
$A$-modules de type fini~; compte tenu de
(\hyperref[I.10.10.2.1-fr]{10.10.2.1}),
(\hyperref[I.10.10.2.3-fr]{10.10.2.3}) et
(\hyperref[I.1.3.12-fr]{1.3.12, (ii)}), on est ramené à montrer que
l'isomorphisme canonique
$\operatorname{Hom}_A(M,N)\xrightarrow{\sim}
\varprojlim_n\operatorname{Hom}_{A_n}(M_n,N_n)$ (avec
$M_n=M\otimes_A A_n$, $N_n=N\otimes_A A_n$) est continu, ce qui a été
prouvé dans (\hyperref[0.7.8.2-fr]{0, 7.8.2}).

\begin{env}[10.11.8]
\label{I.10.11.8-fr}
Comme $\operatorname{Hom}_{\mathscr{O}_{\mathfrak{X}}}
(\mathscr{F},\mathscr{G})$ est le groupe des sections du faisceau de
groupes topologiques
$\mathscr{H}\!om_{\mathscr{O}_{\mathfrak{X}}}
(\mathscr{F},\mathscr{G})$, il est muni d'une topologie de groupe. Si
$\mathfrak{X}$ est \emph{noethérien}, il résulte de
(\hyperref[I.10.11.7.2-fr]{10.11.7.2}) qu'un système fondamental de
voisinages de $0$ dans ce groupe s'obtient en prenant les sous-groupes
$\operatorname{Hom}_{\mathscr{O}_{\mathfrak{X}}}
(\mathscr{F},\mathscr{J}^n\mathscr{G})$ ($n$ arbitraire).
\end{env}

\begin{proposition}[10.11.9]
\label{I.10.11.9-fr}
Soient $\mathfrak{X}$ un préschéma formel noethérien, $\mathscr{F}$ et
$\mathscr{G}$ deux $\mathscr{O}_{\mathfrak{X}}$-Modules cohérents. Dans
le groupe topologique
$\operatorname{Hom}_{\mathscr{O}_{\mathfrak{X}}}
(\mathscr{F},\mathscr{G})$ les homomorphismes surjectifs (resp.
injectifs, bijectifs) forment une partie ouverte.
\end{proposition}

En vertu de (\hyperref[I.10.11.5-fr]{10.11.5}), l'ensemble des
homomorphismes surjectifs dans
$\operatorname{Hom}_{\mathscr{O}_{\mathfrak{X}}}
(\mathscr{F},\mathscr{G})$ est l'image réciproque par l'application
continue
$\operatorname{Hom}_{\mathscr{O}_{\mathfrak{X}}}
(\mathscr{F},\mathscr{G})\to
\operatorname{Hom}_{\mathscr{O}_{X_0}}(\mathscr{F}_0,\mathscr{G}_0)$
d'une partie du groupe discret
$\operatorname{Hom}_{\mathscr{O}_{X_0}}(\mathscr{F}_0,\mathscr{G}_0)$
d'où la première assertion. Pour démontrer la seconde, recouvrons
$\mathfrak{X}$ par un nombre fini d'ouverts formels affines noethériens
$U_i$. Pour que
$\theta\in\operatorname{Hom}_{\mathscr{O}_{\mathfrak{X}}}
(\mathscr{F},\mathscr{G})$ soit injectif, il faut et il suffit que toutes
ses images par les applications (continues) de restriction
$\operatorname{Hom}_{\mathscr{O}_{\mathfrak{X}}}
(\mathscr{F},\mathscr{G})\to
\operatorname{Hom}_{\mathscr{O}_{\mathfrak{X}}|U_i}
(\mathscr{F}|U_i,\mathscr{G}|U_i)$ le soient~; on est donc ramené au cas
affine, et alors cela a déjà été prouvé dans
(\hyperref[0.7.8.3-fr]{0, 7.8.3}).

